\section{Hyperbolic PG-GAT}
\label{sec_meth_hyperbolic}

This section describes a hyperbolic variant of PG-GAT in which the
Euclidean output space $\mathbb{R}^{d}$ is replaced by a Lorentz-model
hyperbolic manifold $\mathbb{H}^{d}$ with curvature $-1$.
The remainder of the PG-GAT architecture is retained,
allowing the experiment to isolate the effect of changing the geometry
of the learned representation.

The motivation follows from the structure of the data.
A human skeleton forms a shallow kinematic tree,
and a multi-person scene can be viewed as a forest of such trees.
Hyperbolic spaces are well suited to representing hierarchical
and tree-structured data because their volume grows exponentially
with radius,
matching the exponential growth of a tree more closely
than the polynomial volume growth of Euclidean space.
Consequently,
tree structures can be represented with lower distortion
in low-dimensional hyperbolic spaces than in comparable
Euclidean spaces~\cite{sarkar2011treeembed, sala2018learningmrep}.

This motivates a direct experimental question for pose grouping.
If part of the representational capacity of Euclidean PG-GAT
is used to encode the hierarchical relationships between joints
in the skeleton,
then placing the final embeddings on a hyperbolic manifold
may provide a more natural geometry for those relationships.
The hyperbolic variant therefore changes the embedding space
while leaving the skeleton-aware message-passing machinery unchanged,
providing a controlled test of whether the geometric prior itself
improves grouping quality or embedding efficiency.

\subsection{Architecture}
\label{sec_meth_hyperbolic_arch}

The internal message-passing layers of the hyperbolic variant remain
Euclidean and retain all three PG-GAT modifications introduced in
Section~\ref{sec_meth_pggat}.
Only the final embedding geometry is changed.
The Euclidean output of the PG-GAT stack is mapped onto the
Lorentz hyperboloid
\begin{equation}
    \mathbb{H}^{d}
    =
    \left\{
        \mathbf{x}\in\mathbb{R}^{d+1}
        :
        \langle\mathbf{x},\mathbf{x}\rangle_{\mathcal{L}}=-1,\;
        x_0>0
    \right\},
    \label{eq:hyperbolic_manifold}
\end{equation}
where the Minkowski inner product is
\begin{equation}
    \langle\mathbf{x},\mathbf{y}\rangle_{\mathcal{L}}
    =
    -x_0y_0+\sum_{k=1}^{d}x_ky_k.
\end{equation}
This hybrid Euclidean-internal, hyperbolic-output design is deliberate.
It isolates the effect of the embedding geometry
while leaving the message-passing architecture unchanged.
A fully hyperbolic PG-GAT would instead require replacing
the linear transformations and aggregation operations
throughout the network with manifold-aware counterparts,
such as M\"obius operations and Fr\'echet-mean aggregation.
Such a model would test several architectural changes at once,
and would therefore make it difficult to attribute any difference
in grouping performance specifically to the output geometry.
The Lorentz model is used instead of the Poincar\'e ball
because of its greater numerical stability
for learned hyperbolic representations~\cite{nickel2018lorentz}.
The implementation uses the \texttt{geoopt} library~\cite{kochurov2020geoopt}.

Before projection onto the manifold,
the Euclidean output $\mathbf{v}_i\in\mathbb{R}^{d}$
is multiplied by a learnable positive tangent-scale parameter $\tau$.
The scaled vector is then mapped from the tangent space at the origin
onto the Lorentz hyperboloid using the exponential map,
\begin{equation}
    \mathbf{x}^{\mathbb{H}}_i
    =
    \mathrm{expmap}_{0}
    \!\left(
        \tau\mathbf{v}_i
    \right).
    \label{eq:hyperbolic_expmap}
\end{equation}
The scale is initialised as
\begin{equation}
    \tau_0=\frac{1}{\sqrt{d}}.
\end{equation}
This initialisation controls the magnitude of the tangent vectors before
the exponential map.
Preliminary measurements showed that it places initial pairwise
geodesic distances approximately in the range $[0.6,1.0]$
for the embedding dimensions considered here.
Without this scaling,
the Euclidean output norm is approximately proportional to $\sqrt{d}$,
and exponential mapping produces initial pairwise hyperbolic distances
on the order of $20$.
At that scale,
the contrastive margin $\gamma=0.5$ provides little useful separation
pressure because most negative pairs already lie far beyond the margin.
The tangent scaling therefore keeps the initial hyperbolic geometry
in a range where the contrastive objective supplies a meaningful gradient.

\subsection{Hyperbolic contrastive loss}
\label{sec_meth_hyperbolic_loss}

The Euclidean contrastive objective of
Equation~\eqref{eq:contrastive_loss} is defined in terms of cosine
similarity between L2-normalised embeddings.
Once the output is moved onto the Lorentz manifold,
cosine similarity is no longer the natural metric.
The hyperbolic variant therefore replaces it with the Lorentz
geodesic distance
\begin{equation}
    d_{\mathbb{H}}(\mathbf{x},\mathbf{y})
    =
    \operatorname{arccosh}
    \left(
        -\langle\mathbf{x},\mathbf{y}\rangle_{\mathcal{L}}
    \right),
    \label{eq:hyperbolic_distance}
\end{equation}
which is non-negative and equals zero when the two points coincide.
Using the same positive and negative pair sets
$\mathcal{P}^{+}$ and $\mathcal{P}^{-}$ defined in
Section~\ref{sec_meth_baseline_loss},
the hyperbolic contrastive objective is
\begin{equation}
    \mathcal{L}_{\mathbb{H}}
    =
    \frac{1}{|\mathcal{P}^{+}|}
    \sum_{(i,j)\in\mathcal{P}^{+}}
    d_{\mathbb{H}}
    \left(
        \mathbf{x}^{\mathbb{H}}_i,
        \mathbf{x}^{\mathbb{H}}_j
    \right)
    +
    \frac{1}{|\mathcal{P}^{-}|}
    \sum_{(i,j)\in\mathcal{P}^{-}}
    \max\!\left(
        0,\,
        \gamma_{\mathbb{H}}
        -
        d_{\mathbb{H}}
        \left(
            \mathbf{x}^{\mathbb{H}}_i,
            \mathbf{x}^{\mathbb{H}}_j
        \right)
    \right).
    \label{eq:hyperbolic_contrastive}
\end{equation}
The first term pulls embeddings belonging to the same person toward
zero geodesic distance.
The second pushes embeddings belonging to different people apart
until their distance exceeds the hyperbolic margin $\gamma_{\mathbb{H}}$.
The margin is set to $\gamma_{\mathbb{H}}=2.0$,
approximately twice the initial mean pairwise distance
produced by the tangent-scale initialisation of
Equation~\eqref{eq:hyperbolic_expmap}.
A separate margin is required because the Euclidean value $\gamma=0.5$
is defined on cosine similarity,
and therefore has no equivalent interpretation
on the Lorentz distance scale.

Computing the full pairwise distance matrix requires an additional
numerical safeguard.
Floating-point error in the batched Minkowski inner products
can move the argument of $\operatorname{arccosh}$
slightly outside its valid domain,
most notably for self-distances whose theoretical argument is exactly one.
The implementation ensures that tensors produced through
\texttt{.expand()} are made contiguous before the pairwise operations,
and applies \texttt{torch.nan\_to\_num} to the resulting distance
matrix as a final safeguard against non-finite values.
These operations do not change the intended loss -
they prevent numerical artefacts in the Lorentz-distance calculation
from propagating through training.

\subsection{Grouping read-out}
\label{sec_meth_hyperbolic_readout}

The downstream $k$-means read-out of
Section~\ref{sec_meth_baseline_training} operates in Euclidean space,
and therefore cannot be applied directly to points on the Lorentz
manifold.
Before clustering,
each hyperbolic embedding is mapped back
to the tangent space at the origin using the logarithmic map,
\begin{equation}
    \mathbf{v}^{\text{tan}}_i
    =
    \mathrm{logmap}_{0}
    \left(
        \mathbf{x}^{\mathbb{H}}_i
    \right),
    \qquad
    \mathrm{logmap}_{0} :
    \mathbb{H}^{d}
    \rightarrow
    T_{\mathbf{0}}\mathbb{H}^{d}.
    \label{eq:hyperbolic_logmap}
\end{equation}
The tangent representation contains the additional Lorentz time
coordinate introduced by the hyperboloid model.
This coordinate is discarded,
leaving a $d$-dimensional spatial vector,
which is then L2-normalised before being supplied to the same $k$-means
procedure used by Euclidean PG-GAT.

The tangent space provides a locally Euclidean approximation
to the curved manifold.
Geometrically,
it can be understood as the flat space that touches the manifold
at the chosen reference point and locally approximates its geometry.
Mapping the embeddings to the tangent space therefore permits
standard Euclidean clustering
while retaining, during training,
the geodesic structure imposed by the hyperbolic contrastive loss.

This approximation is most faithful near the origin.
The tangent-scale parameter $\tau$ introduced in
Equation~\eqref{eq:hyperbolic_expmap} keeps the learned representations
within a comparatively small region of the manifold,
where distortion introduced by the tangent-space projection
is correspondingly limited.
The procedure nevertheless constitutes a deliberate compromise -
training and pairwise supervision use the Lorentz geometry,
whereas the final partition is recovered from its Euclidean
tangent-space approximation.
This makes the hyperbolic variant directly compatible
with the standard PG-GAT $k$-means read-out,
while avoiding the additional methodological change
of introducing a manifold-specific clustering algorithm.

\subsection{Two pre-registered hypotheses and the dimension axis}
\label{sec_meth_hyperbolic_hypotheses}

The hyperbolic variant is evaluated against two hypotheses specified
before the corresponding experiments were run.
The first tests for a quality advantage at matched representational
capacity:

\begin{description}
    \item[H1:] A hyperbolic PG-GAT embedding matches or exceeds the
    grouping performance of its Euclidean counterpart at the same
    output dimension.
\end{description}

The second tests whether the geometric prior permits a more compact
representation:

\begin{description}
    \item[H2:] A hyperbolic PG-GAT embedding matches the grouping
    performance of the Euclidean model at a substantially smaller
    output dimension.
\end{description}

The distinction is important because hyperbolic geometry need not
improve the maximum achievable PGA to be useful.
A model that preserves Euclidean grouping quality
with a substantially lower-dimensional embedding
would constitute an efficiency gain even if H1 is not supported.

The initial evaluation uses two controlled single-run ablations.
H1 is tested at output dimension $d=128$,
comparing Euclidean and hyperbolic embeddings at matched dimensionality.
H2 is tested at $d=32$,
a four-fold reduction relative to the $128$-dimensional comparison
configuration
and an eight-fold reduction relative to the $d=256$ sweep-selected
PG-GAT architecture used for the headline results.
These experiments provide the direct tests of the two pre-specified
hypotheses.

A subsequent eight-run Bayesian sweep over output dimension,
number of message-passing layers,
and hidden dimension
tests whether the single-run conclusions persist across a broader region
of the hyperbolic architecture space.
The sweep is therefore used as defensibility evidence
rather than as the source of the hypotheses -
the hypotheses and their primary comparisons are fixed independently
of the sweep outcome.

Finally, a separate Option~B experiment fine-tunes the
$d=32$ hyperbolic model on COCO.
This experiment tests whether any low-dimensional efficiency advantage
observed after synthetic pre-training survives adaptation
to the real-image distribution,
rather than being specific to the synthetic evaluation regime.

Unless explicitly stated otherwise,
the hyperbolic experiments inherit the PG-GAT architecture components,
training procedure,
optimiser,
and evaluation protocol defined in
Sections~\ref{sec_meth_baseline_training} and
\ref{sec_meth_pggat_training}.
The embedding geometry,
its associated contrastive distance,
and the output dimension under test are therefore the principal
experimental variables.